% บทที่ 2 แนวคิดและงานที่เกี่ยวข้อง
% ย้ายมาจากเค้าโครง: ../chapter/06-literature.tex
\chapter{แนวคิดและงานที่เกี่ยวข้อง}

\section{ทฤษฎีที่เกี่ยวข้อง}

\subsection{สุขภาวะทางจิตและแบบวัด DASS-21}

โครงงานนี้ไม่ได้มองสุขภาวะทางจิตเป็นค่าเดียว แต่แยกเป็นสามด้านที่มีลักษณะต่างกัน
ภาวะซึมเศร้าเน้นอารมณ์ด้านลบแบบหมดพลัง เช่น หมดความสนใจ รู้สึกไร้ค่า มองไม่เห็นอนาคต
ความวิตกกังวลเน้นการตอบสนองของร่างกายต่อสิ่งที่ยังไม่เกิด เช่น ใจสั่น มือสั่น กลัวว่าจะควบคุมตัวเองไม่ได้
ส่วนความเครียดเน้นภาวะตึงตัวเรื้อรัง เช่น หงุดหงิดง่าย ผ่อนคลายยาก ทนรออะไรไม่ได้

การแยกสามด้านนี้มาจากโมเดลสามองค์ประกอบ (tripartite model) ของ Clark และ Watson
\cite{clark1991tripartite} โมเดลนี้อธิบายว่าภาวะซึมเศร้ากับความวิตกกังวลมีส่วนที่ทับซ้อนกัน
คือความทุกข์ทั่วไป และมีส่วนที่แยกจากกันชัดเจน ภาวะซึมเศร้ามาคู่กับการขาดอารมณ์ด้านบวก
ส่วนความวิตกกังวลมาคู่กับการตื่นตัวทางกาย
ถ้าใช้แบบวัดที่ให้คะแนนรวมค่าเดียว ก็จะตอบไม่ได้ว่าการใช้โซเชียลมีเดียสัมพันธ์กับด้านใดมากกว่ากัน

แบบวัด DASS-21 พัฒนามาจากคู่มือของ Lovibond และ Lovibond \cite{lovibond1995manual}
เป็นแบบวัดเชิงอาการ (symptom-based) ที่ถามถึงความรู้สึกในช่วง 1 สัปดาห์ที่ผ่านมา
มี 21 ข้อ แบ่งเป็น 3 มิติ มิติละ 7 ข้อ แต่ละข้อให้คะแนน 0 ถึง 3
รวมคะแนนรายมิติแล้วคูณ 2 เพื่อเทียบกับเกณฑ์ของฉบับเต็ม 42 ข้อ
คะแนนยิ่งสูงแปลว่าอาการยิ่งมาก DASS-21 จึงเป็นแบบวัดเชิงลบ
ต่างจากแบบวัดสุขภาวะเชิงบวกที่คะแนนสูงแปลว่าดี เรื่องนี้ต้องระวังตอนตีความเครื่องหมายของค่าสัมประสิทธิ์ในบทที่ 4

แบบวัดใช้ได้ก็ต่อเมื่อผ่านการตรวจสอบคุณภาพในกลุ่มที่ใกล้เคียงกับผู้ตอบจริง
คุณสมบัติที่มักรายงานมีสองส่วน ส่วนแรกคือความเชื่อมั่นภายใน (internal consistency)
ซึ่งดูว่าข้อคำถามในมิติเดียวกันวัดเรื่องเดียวกันจริงหรือไม่ นิยมรายงานด้วยค่าสัมประสิทธิ์แอลฟาของครอนบาค
และถือว่าค่าตั้งแต่ 0.70 ขึ้นไปใช้ได้ ส่วนที่สองคือโครงสร้างองค์ประกอบ (factor structure)
ซึ่งตรวจว่าการแบ่งข้อคำถามออกเป็นสามมิติเข้ากับข้อมูลจริงหรือไม่
DASS-21 มีฉบับแปลภาษาไทยเผยแพร่บนเว็บไซต์ทางการของแบบวัด \cite{webster2008thai}
และมีการตรวจสอบคุณภาพในกลุ่มนักศึกษาไทยแล้ว \cite{wittayapun2023validation}
โครงงานนี้จึงใช้ฉบับแปลดังกล่าวโดยไม่แก้ข้อความ
และคำนวณค่าแอลฟาของทั้งสามมิติจากข้อมูลของโครงงานเองอีกครั้ง (หัวข้อ~\ref{sec:reliability})

\subsection{พฤติกรรมการใช้โซเชียลมีเดียและ doomscrolling}

งานที่ศึกษาผลของโซเชียลมีเดียต่อสุขภาพจิตมักแยกตัวแปรเป็นสองแบบ
แบบแรกคือ\textit{ปริมาณ} หรือเวลาที่ใช้ต่อวัน แบบที่สองคือ\textit{ลักษณะการใช้}
ซึ่งรวมประเภทเนื้อหาที่เสพและรูปแบบการมีปฏิสัมพันธ์
การวัดแค่ปริมาณมีข้อจำกัด เพราะคนที่ใช้เวลาเท่ากันแต่ดูเนื้อหาต่างกันอาจได้รับผลต่างกันมาก
โครงงานนี้จึงเก็บทั้งชั่วโมงการใช้งาน ประเภทเนื้อหาที่ดูมากที่สุด และจำนวนวันที่ติดตามข่าวเชิงลบ

กลไกที่มักใช้อธิบายความสัมพันธ์ระหว่างโซเชียลมีเดียกับสุขภาวะทางจิตมีสามทาง

\begin{enumerate}[label=\arabic*),leftmargin=1.9cm,itemsep=0.15em,topsep=0.2em]
    \item \textbf{การแย่งเวลา (displacement)} เวลาหน้าจอไปเบียดเวลานอน เวลาออกกำลังกาย
          และเวลาพบปะผู้คนจริง ซึ่งล้วนเกี่ยวข้องกับสุขภาพจิต
          กลไกนี้เป็นเหตุผลที่โครงงานเก็บชั่วโมงการนอน คุณภาพการนอน และจำนวนวันที่ออกกำลังกายเป็นตัวแปรควบคุม
    \item \textbf{การเปรียบเทียบทางสังคม (social comparison)} ทฤษฎีของ Festinger
          \cite{festinger1954social} อธิบายว่าคนเรามักประเมินตัวเองโดยเทียบกับผู้อื่น
          เนื้อหาบนโซเชียลมีเดียส่วนใหญ่เป็นด้านที่คัดมาแล้ว
          การเทียบตัวเองกับภาพเหล่านั้นจึงมักเป็นการเทียบขึ้น (upward comparison)
          ซึ่งสัมพันธ์กับความรู้สึกด้อยค่าและอารมณ์ด้านลบ
          กลไกนี้เป็นที่มาของตัวเลือกเนื้อหาประเภท ``ชีวิตประจำวันของคนอื่น'' ในแบบสอบถาม
    \item \textbf{การเสพเนื้อหาด้านลบต่อเนื่อง (doomscrolling)} คือการไถดูข่าวหรือเนื้อหาเชิงลบ
          ไปเรื่อย~ๆ จนหยุดยาก ทั้งที่รู้ว่ายิ่งดูยิ่งรู้สึกแย่
          คำนี้เริ่มใช้แพร่หลายช่วงการระบาดของโควิด-19
          และมีผู้อธิบายว่าเป็นความพยายามลดความไม่แน่นอนด้วยการตามข่าวตลอดเวลา
          ซึ่งกลับกลายเป็นวงจรที่เพิ่มความตึงเครียด
          อคติเชิงลบของสื่อและระบบแนะนำเนื้อหาของแพลตฟอร์มยิ่งเร่งวงจรนี้
\end{enumerate}

\noindent
doomscrolling มีแบบวัดเฉพาะ เช่น Doomscrolling Scale
แต่ถ้าเพิ่มแบบวัดนี้เข้าไป แบบสอบถามจะยาวเกินไป
โครงงานนี้จึงใช้ตัวชี้วัดอย่างง่ายแทน คือจำนวนวันใน 7 วันที่ผ่านมาที่ผู้ตอบติดตามข่าวเชิงลบ
ตัวชี้วัดนี้วัดความถี่ของพฤติกรรม ไม่ได้วัดลักษณะการติดข่าว
บทที่ 5 อภิปรายข้อจำกัดนี้

\subsection{การวิเคราะห์สหสัมพันธ์และการถดถอยพหุคูณ}

\textbf{สหสัมพันธ์เพียร์สัน (Pearson correlation)} วัดความสัมพันธ์เชิงเส้นระหว่างตัวแปรเชิงปริมาณสองตัว
ค่า $r$ อยู่ในช่วง $-1$ ถึง $1$ เครื่องหมายบอกทิศทาง ขนาดบอกความแรง
เมื่อยกกำลังสองจะได้ $r^2$ ซึ่งตีความได้ว่าตัวแปรหนึ่งอธิบายความแปรปรวนของอีกตัวได้กี่เปอร์เซ็นต์
เกณฑ์ที่นิยมใช้ตีความขนาดของ $r$ คือเกณฑ์ของ Cohen \cite{cohen1988statistical}
คือ 0.10 เล็ก 0.30 ปานกลาง และ 0.50 ใหญ่
วิธีนี้มีข้อควรระวังสองข้อ ข้อแรกคือ $r$ จับได้เฉพาะความสัมพันธ์เชิงเส้น
ถ้าความสัมพันธ์เป็นรูปตัวยู ค่าที่ได้จะใกล้ศูนย์ทั้งที่มีความสัมพันธ์อยู่ จึงต้องดู scatter plot ประกอบทุกครั้ง
ข้อที่สองคือความสัมพันธ์ไม่ได้แปลว่าเป็นเหตุเป็นผล
ข้อมูลแบบตัดขวาง (cross-sectional) เก็บทุกตัวแปรพร้อมกันครั้งเดียว จึงบอกลำดับก่อนหลังไม่ได้

\textbf{การถดถอยเชิงเส้นพหุคูณ (multiple linear regression)} ใช้ดูผลของตัวแปรต้นหลายตัวพร้อมกัน
สมการมีรูปแบบดังนี้

\begin{center}
$Y = \beta_0 + \beta_1 X_1 + \beta_2 X_2 + \cdots + \beta_k X_k + \varepsilon$
\end{center}

\noindent
$\beta_i$ บอกว่าเมื่อ $X_i$ เพิ่มขึ้น 1 หน่วย ค่า $Y$ เปลี่ยนไปเท่าใด
\textit{เมื่อตัวแปรต้นตัวอื่นคงที่} จุดนี้ทำให้การถดถอยต่างจากสหสัมพันธ์
เพราะแยกผลของโซเชียลมีเดียออกจากผลของการนอนและความกดดันจากการเรียนได้
ที่ต้องแยกเพราะมีปัญหาตัวแปรกวน (confounding variable)
เช่น ถ้าพบว่าชั่วโมงการใช้โซเชียลมีเดียสัมพันธ์กับคะแนนความเครียด
ก็ยังเป็นไปได้ว่าคนที่นอนน้อยทั้งเล่นมือถือมากและเครียดมาก
การใส่ตัวแปรการนอนเข้าไปในสมการช่วยตรวจว่าความสัมพันธ์เดิมยังเหลืออยู่หรือไม่

ค่าที่ต้องดูประกอบกันมีสามค่า ค่า $R^2$ บอกว่าสมการอธิบายความแปรปรวนของตัวแปรตามได้กี่เปอร์เซ็นต์
ค่า $p$ ของแต่ละสัมประสิทธิ์บอกว่าผลนั้นน่าจะเกิดจากความบังเอิญหรือไม่
ส่วนค่า VIF (variance inflation factor) ใช้ตรวจว่าตัวแปรต้นสัมพันธ์กันเองมากเกินไปหรือไม่ (multicollinearity)
ปัญหานี้ทำให้ค่าสัมประสิทธิ์แกว่งและตีความผิดได้ โดยทั่วไปถือว่า VIF เกิน 5 เริ่มน่ากังวล
การถดถอยเชิงเส้นยังมีข้อตกลงเบื้องต้นที่ต้องตรวจ ได้แก่ ความสัมพันธ์เป็นเชิงเส้น
ค่าคลาดเคลื่อนแจกแจงแบบปกติ ความแปรปรวนของค่าคลาดเคลื่อนคงที่ และค่าคลาดเคลื่อนเป็นอิสระต่อกัน
ตรวจได้ด้วย residual plot และ Q-Q plot

\subsection{สหสัมพันธ์บางส่วนและการเปรียบเทียบค่าเฉลี่ยระหว่างกลุ่ม}

\textbf{สหสัมพันธ์บางส่วน (partial correlation)} คือค่า $r$ ระหว่างตัวแปรสองตัว
หลังหักส่วนที่ตัวแปรควบคุมอธิบายได้ออกจากทั้งสองตัวแล้ว
วิธีคำนวณคือถดถอยตัวแปรแต่ละตัวด้วยตัวแปรควบคุม แล้วหาสหสัมพันธ์ระหว่างค่าคลาดเคลื่อนที่เหลือ
ค่าที่ได้ตอบคำถามเดียวกับการถดถอยพหุคูณ คือความสัมพันธ์ยังเหลืออยู่หรือไม่เมื่อตัวแปรควบคุมเท่ากัน
แต่รายงานในหน่วยเดียวกับ $r$ จึงเทียบกับสหสัมพันธ์ปกติได้ตรง~ๆ

\textbf{การทดสอบที (t-test)} ใช้เทียบค่าเฉลี่ยของกลุ่มเดียวกับค่าที่กำหนด (one-sample)
หรือเทียบค่าเฉลี่ยของสองกลุ่มที่เป็นอิสระต่อกัน (independent samples)
ถ้ามีมากกว่าสองกลุ่มจะใช้\textbf{การวิเคราะห์ความแปรปรวนทางเดียว (one-way ANOVA)}
ซึ่งทดสอบว่ามีอย่างน้อยหนึ่งกลุ่มที่ค่าเฉลี่ยต่างจากกลุ่มอื่นหรือไม่
ขนาดผลของ ANOVA รายงานด้วย $\eta^2$ คือสัดส่วนความแปรปรวนของตัวแปรตามที่อธิบายได้ด้วยการแบ่งกลุ่ม
ทั้งสองวิธีมีข้อตกลงเบื้องต้นว่าข้อมูลในแต่ละกลุ่มแจกแจงใกล้ปกติและความแปรปรวนใกล้เคียงกัน
ถ้าข้อตกลงนี้ไม่น่าเชื่อถือ เช่น กลุ่มเล็กหรือคะแนนเบ้มาก
จะใช้\textbf{การทดสอบ Kruskal-Wallis} ซึ่งเทียบลำดับที่ของข้อมูลแทนค่าจริง เป็นการตรวจซ้ำ

\textbf{การปรับค่า $p$ ด้วยวิธี Bonferroni} ใช้เมื่อทดสอบหลายครั้งพร้อมกัน
ยิ่งทดสอบมากครั้ง ยิ่งมีโอกาสได้ผลที่มีนัยสำคัญโดยบังเอิญ
วิธีนี้คูณค่า $p$ ด้วยจำนวนการทดสอบ หรือเทียบได้กับการลดระดับนัยสำคัญจาก 0.05 เป็น 0.05 หารด้วยจำนวนการทดสอบ

\section{งานวิจัยที่เกี่ยวข้อง}

\subsection{การใช้โซเชียลมีเดียกับสุขภาพจิตของวัยรุ่น}

Sala, Porcaro และ Gómez \cite{sala2024umbrella} ทำการทบทวนวรรณกรรมแบบร่ม (umbrella review)
คือการทบทวนงานที่เป็นการทบทวนวรรณกรรมและการวิเคราะห์อภิมานอีกชั้นหนึ่ง
คณะผู้วิจัยค้นจากฐานข้อมูล Scopus, Web of Science, PsycInfo และ PubMed
ครอบคลุมงานที่ตีพิมพ์ระหว่างเดือนมกราคม 2015 ถึงเดือนเมษายน 2023
คัดกรองจาก 1,470 รายการ เหลืองานที่ผ่านการประเมินคุณภาพด้วยเกณฑ์ AMSTAR 2 จำนวน 24 รายการ

ผลการทบทวนพบว่าการใช้โซเชียลมีเดียมีทั้งความเสี่ยงและโอกาสต่อสุขภาพจิตของวัยรุ่น
ผลที่เกิดขึ้นไม่ได้ขึ้นกับเวลาที่ใช้อย่างเดียว
แต่ขึ้นกับลักษณะของผู้ใช้ รูปแบบการใช้งาน และการออกแบบของแพลตฟอร์มด้วย
ผู้วิจัยยังชี้ว่างานที่ผ่านมาส่วนใหญ่มองวัยรุ่นและแพลตฟอร์มเป็นก้อนเดียวกัน
ทั้งที่หลากหลายมาก จึงเสนอให้ศึกษาแยกตามกลุ่มผู้ใช้และรูปแบบการใช้งานมากขึ้น

งานนี้เกี่ยวข้องกับโครงงานสองทาง ทางแรกคือเป็นฐานเปรียบเทียบว่าความสัมพันธ์ที่พบในงานระดับสากลมีขนาดเล็ก
จึงไม่ควรคาดหวังค่า $r$ ที่สูงจากโครงงานนี้
ทางที่สองคือสนับสนุนการออกแบบที่เก็บทั้งชั่วโมงการใช้งานและประเภทเนื้อหา
แทนที่จะวัดแค่เวลาที่ใช้

\subsection{การเสพข่าวเชิงลบต่อเนื่องกับความวิตกกังวลและภาวะซึมเศร้า}

Dominguez-Rodriguez และคณะ \cite{dominguez2025climate} ศึกษาความสัมพันธ์ระหว่าง doomscrolling
แบบทั่วไปกับ doomscrolling เฉพาะข่าวการเปลี่ยนแปลงสภาพภูมิอากาศ
งานนี้เก็บข้อมูลแบบตัดขวางด้วยแบบสอบถามออนไลน์จากผู้ตอบในเนเธอร์แลนด์และเยอรมนี
และวิเคราะห์จากกลุ่มตัวอย่างสุดท้าย 365 คน
แบบวัดที่ใช้คือ Doomscrolling Scale (DS) 15 ข้อ และ Climate Change Doomscrolling Scale (CCDS) 11 ข้อ
ที่พัฒนาขึ้นใหม่ในงานนี้ ปัจจัยเสี่ยงวัดด้วยแบบวัดภาวะซึมเศร้าและความวิตกกังวล
ปัจจัยปกป้องวัดด้วยแบบวัดการสนับสนุนทางสังคมและทักษะการรับมือปัญหา
สถิติที่ใช้คือสหสัมพันธ์เพียร์สันและการถดถอยเชิงเส้นพหุคูณ โดยตรวจผลของเพศในฐานะตัวแปรกำกับด้วย

doomscrolling แบบทั่วไปสัมพันธ์ทางบวกกับ doomscrolling เรื่องสภาพภูมิอากาศ
อย่างมีนัยสำคัญ ($r = 0.35$, $p < 0.001$) ซึ่งเป็นความสัมพันธ์ขนาดปานกลาง
ด้านปัจจัยเสี่ยง คะแนนความวิตกกังวลทำนาย doomscrolling ได้อย่างมีนัยสำคัญ ($p < 0.001$)
ส่วนคะแนนภาวะซึมเศร้าไม่มีนัยสำคัญเมื่อดูภาพรวม
แต่เมื่อแยกตามเพศ ทิศทางกลับต่างกัน คือเป็นบวกในผู้หญิงและเป็นลบในผู้ชาย
ด้านปัจจัยปกป้อง การสนับสนุนทางสังคมไม่มีผลอย่างมีนัยสำคัญ
ส่วนทักษะการรับมือปัญหามีผลในทางปกป้อง แต่ขนาดผลไม่มาก

งานนี้สนับสนุนคำถามวิจัยข้อ 3 ของโครงงานโดยตรง ว่าการเสพข่าวเชิงลบน่าจะสัมพันธ์กับความวิตกกังวล
มากกว่ามิติอื่น และเป็นเหตุผลที่โครงงานเลือกแบบวัดที่แยกคะแนนสามมิติ
ข้อแตกต่างที่ต้องระวังคืองานนี้ใช้แบบวัด doomscrolling เต็มรูปแบบ
ขณะที่โครงงานนี้ใช้แค่จำนวนวันที่ติดตามข่าว ความสัมพันธ์ที่ได้จึงอาจมีขนาดเล็กกว่า

\subsection{คุณสมบัติเชิงจิตมิติของแบบวัด DASS-21}

Manzar และคณะ \cite{manzar2025dass21} ตรวจสอบคุณสมบัติเชิงจิตมิติของ DASS-21
ในนักศึกษามหาวิทยาลัยในเอธิโอเปีย 364 คน
โดยใช้ทั้งทฤษฎีการทดสอบแบบดั้งเดิม (classical test theory) และทฤษฎีการตอบสนองข้อสอบ (item response theory)
การวิเคราะห์ประกอบด้วยการวิเคราะห์องค์ประกอบจากเมทริกซ์สหสัมพันธ์แบบโพลีคอริก แบบจำลองราสช์
และการตรวจการทำหน้าที่ต่างกันของข้อคำถามระหว่างเพศ (differential item functioning)

แบบวัดมีความเชื่อมั่นภายในสูงมาก (แอลฟาของครอนบาค $= 0.93$ และโอเมกาของแมคโดนัลด์ $= 0.93$)
และข้อคำถามทำหน้าที่ไม่ต่างกันระหว่างเพศ จึงเทียบคะแนนระหว่างเพศได้
แต่มีข้อคำถาม 3 ข้อที่มีค่าพารามิเตอร์ไม่ดีนัก
และการวิเคราะห์องค์ประกอบในกลุ่มตัวอย่างนี้สนับสนุนโครงสร้างแบบองค์ประกอบเดียวมากกว่าสามองค์ประกอบ
คณะผู้วิจัยสรุปว่าโดยรวมแบบวัดมีความเที่ยงตรงเชิงจิตมิติที่ดีในกลุ่มนักศึกษามหาวิทยาลัย

งานนี้มีนัยต่อโครงงานสองข้อ ข้อแรกคือยืนยันว่า DASS-21 ใช้ได้ดีกับนักศึกษามหาวิทยาลัย
ซึ่งเป็นกลุ่มเดียวกับผู้ตอบของโครงงาน ข้อที่สองคือเตือนว่าโครงสร้างสามมิติอาจไม่ชัดเจนในทุกบริบท
ถ้าคะแนนทั้งสามมิติสัมพันธ์กันเองสูงมาก การเทียบว่า
``โซเชียลมีเดียสัมพันธ์กับมิติใดมากที่สุด'' ต้องตีความอย่างระมัดระวัง

\subsection{ขนาดความสัมพันธ์ในกลุ่มตัวอย่างขนาดใหญ่}

Orben และ Przybylski \cite{orben2019association} วิเคราะห์ข้อมูลวัยรุ่นจากชุดข้อมูลขนาดใหญ่ 3 ชุด
รวม 355,358 คน ด้วยวิธี specification curve analysis
คือคำนวณความสัมพันธ์ซ้ำในทุกวิธีวิเคราะห์ที่สมเหตุสมผล เช่น เลือกตัววัดการใช้เทคโนโลยี ตัววัดสุขภาวะ
และตัวแปรควบคุมต่างกัน แล้วดูผลทั้งชุดแทนการเลือกรายงานผลเพียงแบบเดียว

ผลพบว่าการใช้เทคโนโลยีดิจิทัลสัมพันธ์ทางลบกับสุขภาวะของวัยรุ่น แต่ขนาดเล็กมาก
อธิบายความแปรปรวนของสุขภาวะได้มากที่สุดเพียงร้อยละ 0.4

งานนี้กำหนดกรอบการตีความผลของโครงงาน ความสัมพันธ์ขนาดร้อยละ 0.4 ตรงกับค่า $r$ ประมาณ 0.06
ซึ่งกลุ่มตัวอย่างหลักแสนคนตรวจพบได้ แต่กลุ่มตัวอย่างหลักร้อยคนมีโอกาสตรวจพบต่ำมาก
ถ้าโครงงานนี้ไม่พบความสัมพันธ์ระหว่างเวลาที่ใช้กับคะแนน DASS-21
ผลนั้นจึงสอดคล้องกับงานขนาดใหญ่ ไม่ได้แปลว่าไม่มีความสัมพันธ์เลย

\section{ช่องว่างของงานวิจัยที่โครงงานนี้เข้าไปเติม}

งานทบทวนภาพรวมพบว่าความสัมพันธ์มีขนาดเล็กและขึ้นกับรูปแบบการใช้งาน
และงานที่ใช้กลุ่มตัวอย่างกว่าสามแสนคนยืนยันว่าความสัมพันธ์กับเวลาที่ใช้มีขนาดเล็กมาก
งานเรื่อง doomscrolling ทำในบริบทยุโรปและใช้แบบวัดเฉพาะทาง
ส่วนงานตรวจสอบแบบวัดทำในนักศึกษาแอฟริกา
คณะผู้จัดทำยังไม่พบงานที่ดูพร้อมกันทั้งชั่วโมงการใช้งานที่อ่านจากค่าที่เครื่องบันทึกไว้ ประเภทเนื้อหาที่เสพ
และคะแนนสามมิติของ DASS-21 ในกลุ่มนักศึกษาไทย
โครงงานนี้เติมช่องว่างดังกล่าวในขอบเขตเล็ก~ๆ คือมหาวิทยาลัยเดียว
และบทที่ 5 นำผลที่ได้ไปเทียบกับงานทั้งสี่ชิ้น
